% Chapter introduction: why the topic arose, who first asked, what it answers.
% STATUS: written 2026-11-06 (English gate clean; sources read in full as recorded below).
% Summarises Sections 1--10 of this chapter; every claim is cited in its section.

A derivation of an algebra $A$ is a linear map $D\colon A \to A$ with
\[
  D(xy) = D(x)\,y + x\,D(y).
\]
The definition is short, but it came late. The product rule is as old as the
calculus. For two and a half centuries it was only a property of the
derivative of a function. It became a definition only in the 1920s and 1930s.
Then a derivation became a map on an abstract algebra, fixed by this one
identity.

This chapter tells how that happened. The story has three parts, and
Figure~\ref{fig:origins:timeline} places its sources in time. We use modern
notation throughout. $A$ is an algebra, $M$ a bimodule, $K$ and $k$ are fields,
and $G$ is a group. Maps are written on the left. When a source uses other
notation, we say so once.

\hbfigure{01-origins/timeline}{The sources of Chapter~1 in order of publication.}{fig:origins:timeline}

The first part stays with commutative rings and fields. It starts in 1684,
when Leibniz states the product rule for his differences
(Section~\ref{sec:origins:leibniz-rule}). Two centuries later, Dedekind and
Weber need a derivative of algebraic functions without any limits. In 1882
they define it by algebra alone, and in 1897 Hilbert writes the different of
an algebraic number as a derivative (Section~\ref{sec:origins:function-fields}).

In characteristic $p$ there is a new difficulty. A nonzero polynomial can
have zero derivative; for example, $x^p$ does
\cite[\S 9, p.~213]{Steinitz1910}. Steinitz
meets this in 1910. Later it leads Jacobson to a Galois theory in which
derivations take the place of automorphisms
(Section~\ref{sec:origins:fields-and-separability}). It also leads Hasse and
Schmidt to higher derivations (Section~\ref{sec:origins:higher-derivations}).
Meanwhile Ritt and his students study fields with a derivation. For them such
a field is the natural setting for differential equations
(Section~\ref{sec:origins:differential-algebra}).

The second part leaves commutativity behind. In 1925 Born and Jordan keep the
product rule for matrices. In the same year Dirac shows that every such
operation on quantum variables is a commutator
(Section~\ref{sec:origins:physics}). In 1933 Ore builds polynomials over a
skew field. He finds a product rule twisted by an endomorphism
(Section~\ref{sec:origins:ore-and-twists}). In 1937 Jacobson studies the
derivations of any algebra. He shows that they form a Lie algebra, and in
characteristic $p$ a restricted one (Section~\ref{sec:origins:lie-theory}).

The third part shows how these lines meet. Since Hilbert, number theorists
had studied maps called crossed homomorphisms. They come from groups that act
on fields (Section~\ref{sec:origins:crossed-homomorphisms}). In 1945
Hochschild puts the derivations of an algebra into a cohomology theory.
Derivations are the one-cocycles, and inner derivations are the coboundaries
(Section~\ref{sec:origins:cohomology}). After this step the derivation is a
central object of modern algebra.

Along the way, three facts appear again and again. They are also the first
questions of the rest of the book.

The first fact concerns constants. A derivation is fixed by one identity, and
its kernel is the set of constants. This kernel carries information about the
algebra. Automorphisms have fixed points in the same way. Dedekind, Weber,
Steinitz and Jacobson all use this link
(Sections~\ref{sec:origins:function-fields}
and~\ref{sec:origins:fields-and-separability}).

The second fact concerns brackets. The commutator of two derivations is again
a derivation. So the derivations of any algebra form a Lie algebra. Dirac
meets this fact in physics, and Jacobson makes it the base of a theory
(Sections~\ref{sec:origins:physics} and~\ref{sec:origins:lie-theory}).

The third fact concerns inner derivations. For every element $a$, the map
$x \mapsto ax - xa$ is a derivation. When does every derivation have this
form? Dirac asked this for matrices. Hochschild answered it in general with
the group $H^1$ (Sections~\ref{sec:origins:crossed-homomorphisms}
and~\ref{sec:origins:cohomology}). The question runs through every later
chapter.
